\chapter*{Abstract}
\addcontentsline{toc}{chapter}{Abstract}

Nebula is a simulation-in-the-loop framework developed for the Astera Labs
Nebula Competition problem statement, ``AI/ML for Analog Design.'' It explores the design space of a
PCI Express Generation~2 oriented receiver equalizer. The implemented signal
path combines a parameterized SKY130 source-degenerated continuous-time linear
equalizer (CTLE), a file-driven four-port S-parameter channel, and a behavioral
one-tap decision-feedback equalizer (DFE). Candidate designs can be generated by
Proximal Policy Optimization (PPO), Random Search, or the Cross-Entropy Method
(CEM), then evaluated by a staged pipeline of DC, AC, transient, channel, noise,
and third-harmonic distortion analyses using ngspice.

The project contributes a reusable circuit-evaluation interface, explicit
failure semantics, content-addressed caching, provenance records, versioned PPO
state and reward formulations, checkpoint handling, and optional process,
voltage, and temperature (PVT) qualification. Recorded real-SPICE artifacts show
that the pipeline can produce nominally feasible candidates. Selected design
\code{pipeline\_ep0} passed all 27 conditions in the demonstrated reduced
TT/SS/FF PVT grid at 512-bit candidate fidelity using the public IEEE/IBM
development channel. This result is not PCI-SIG compliance evidence and does
not cover the SF/FS corners in the competition specification.
In the only matched-budget, no-warm-start comparison, Random Search found one
feasible design in 20 trials, whereas PPO and CEM found none. The single-seed
experiment therefore does not demonstrate PPO superiority.

The present work satisfies the requested architectural objective--a target-driven
Python framework that connects reinforcement-learning candidate generation to a
SPICE simulator and emits parameters, metrics, schematics, and run evidence--but
it does not yet establish the competition's stronger efficiency and full
specification-closure goals. It is a Stage~1 design-automation study rather than a complete
physical receiver. The DFE is behavioral, the CTLE contains ideal passive
elements and an ideal tail-current source, and post-layout, jitter, mismatch,
BER, total-area, and total-power signoff remain future work.

\textbf{Keywords:} analog design automation, CTLE, DFE, reinforcement learning,
PPO, ngspice, SKY130, PVT analysis.
